% GENERATED FILE. Edit canonical lessons, then run npm run generate:books.
% canonical-source-hash: fnv1a64:4c31e6847246dce1
% canonical-lessons: UR-C280-daku, UR-C280-daka, UR-C280-banduq, UR-C280-bandhna, UR-C280-lutera

\chapter{Robbers on the Road}
\label{ch:ur-robbers-on-the-road}
\hlchaptermodality{280}

\begin{chapteropening}
\textbf{By the end of this chapter:} \emph{I can tell what robbers did and say what we will do if they come.}

Complete the last lesson of chapter 280: I can tell what robbers did and say what we will do if they come.
\end{chapteropening}

\section[\texorpdfstring{\ur{ڈاکو} (ḍākū)}{ڈاکو (ḍākū)}]{\ur{ڈاکو} (ḍākū) — a bandit}

\label{lesson:UR-C280-daku}



\begin{quote}
Before the new one: say the Urdu for a scene, then the Urdu for to be seen.
\end{quote}

\subsection*{You'll want to know: \ur{ڈاکو}}
\textbf{\ur{ڈاکو}} — \emph{ḍākū} — \textquotedblleft{}a bandit\textquotedblright{}.

\begin{quote}
\ur{رات} \ur{کو} \ur{سڑک} \ur{پر} \ur{ڈاکو} \ur{تھے۔}
\end{quote}

\begin{quote}
\emph{rāt ko saṛak par ḍākū the.}
\end{quote}

\begin{quote}
There were bandits on the road at night.
\end{quote}

Masculine.

\begin{cousinweb}
is inherited: CDIAL lists the Hindi cousin \emph{ḍākū}, \textquotedblleft{}robber\textquotedblright{}, under \emph{*ḍākka}, a form it reconstructs (CDIAL 5543). A \textbf{\ur{ڈاکو}} robs with a gang, on the roads; a \textbf{\ur{چور}} steals alone, in secret.
\end{cousinweb}

\begin{grammarlens}[title={What you've built}]
One more word for telling what happened.
\end{grammarlens}

\subsection*{Your turn}
\begin{itemize}
\raggedright
  \item \emph{Say it:} \emph{ḍākū}
  \item \emph{Say it:} \emph{ḍākū}, once more
  \item \emph{Say it:} \emph{rāt ko saṛak par ḍākū the.} — There were bandits on the road at night.
  \item \emph{Recall:} say the Urdu for a scene, then the Urdu for to be seen, then say \emph{ḍākū} again
\end{itemize}

\begin{tcolorbox}[breakable,colback=teal!4,colframe=teal!35!black,title={Before you move on}]
What does \ur{ڈاکو} mean? (\textquotedblleft{}A bandit\textquotedblright{}.) Say it once more.
\end{tcolorbox}
\section[\texorpdfstring{\ur{ڈاکہ} (ḍākā)}{ڈاکہ (ḍākā)}]{\ur{ڈاکہ} (ḍākā) — an armed robbery}

\label{lesson:UR-C280-daka}



\begin{quote}
Before the new one: say the Urdu for to be seen, then the Urdu for a bandit.
\end{quote}

\subsection*{You'll want to know: \ur{ڈاکہ}}
\textbf{\ur{ڈاکہ}} — \emph{ḍākā} — \textquotedblleft{}an armed robbery\textquotedblright{}.

\begin{quote}
\ur{کل} \ur{بینک} \ur{میں} \ur{ڈاکہ} \ur{پڑا۔}
\end{quote}

\begin{quote}
\emph{kal baink meṅ ḍākā paṛā.}
\end{quote}

\begin{quote}
Yesterday the bank was robbed.
\end{quote}

Masculine.

\begin{cousinweb}
\textbf{\ur{ڈاکہ}} (\emph{ḍākā}) is inherited: CDIAL lists the Hindi cousin \emph{ḍākā}, \textquotedblleft{}robbery, a band of robbers\textquotedblright{}, under \emph{*ḍākka} (CDIAL 5543), the entry that also holds \textbf{\ur{ڈاکو}}. A robbery is said to fall: \textbf{\ur{ڈاکہ} \ur{پڑا}}.
\end{cousinweb}

\begin{grammarlens}[title={What you've built}]
One more word for telling what happened.
\end{grammarlens}

\subsection*{Your turn}
\begin{itemize}
\raggedright
  \item \emph{Say it:} \emph{ḍākā}
  \item \emph{Say it:} \emph{ḍākā}, once more
  \item \emph{Say it:} \emph{kal baink meṅ ḍākā paṛā.} — Yesterday the bank was robbed.
  \item \emph{Recall:} say the Urdu for to be seen, then the Urdu for a bandit, then say \emph{ḍākā} again
\end{itemize}

\begin{tcolorbox}[breakable,colback=teal!4,colframe=teal!35!black,title={Before you move on}]
What does \ur{ڈاکہ} mean? (\textquotedblleft{}An armed robbery\textquotedblright{}.) Say it once more.
\end{tcolorbox}
\section[\texorpdfstring{\ur{بندوق} (bandūq)}{بندوق (bandūq)}]{\ur{بندوق} (bandūq) — a gun}

\label{lesson:UR-C280-banduq}



\begin{quote}
Before the new one: say the Urdu for a bandit, then the Urdu for an armed robbery.
\end{quote}

\subsection*{You'll want to know: \ur{بندوق}}
\textbf{\ur{بندوق}} — \emph{bandūq} — \textquotedblleft{}a gun\textquotedblright{}.

\begin{quote}
\ur{ڈاکو} \ur{کے} \ur{پاس} \ur{بندوق} \ur{تھی۔}
\end{quote}

\begin{quote}
\emph{ḍākū ke pās bandūq thī.}
\end{quote}

\begin{quote}
The bandit had a gun.
\end{quote}

Feminine.

\begin{cousinweb}
\textbf{\ur{بندوق}} (\emph{bandūq}) is usually traced, through Persian, to Arabic \emph{bunduq}, \textquotedblleft{}a hazelnut\textquotedblright{}, and so a pellet; the sense of a firearm is thought to have grown from the pellets it fired. It is feminine: \textbf{\ur{بندوق} \ur{تھی}}.
\end{cousinweb}

\begin{grammarlens}[title={What you've built}]
One more word for telling what happened.
\end{grammarlens}

\subsection*{Your turn}
\begin{itemize}
\raggedright
  \item \emph{Say it:} \emph{bandūq}
  \item \emph{Say it:} \emph{bandūq}, once more
  \item \emph{Say it:} \emph{ḍākū ke pās bandūq thī.} — The bandit had a gun.
  \item \emph{Recall:} say the Urdu for a bandit, then the Urdu for an armed robbery, then say \emph{bandūq} again
\end{itemize}

\begin{tcolorbox}[breakable,colback=teal!4,colframe=teal!35!black,title={Before you move on}]
What does \ur{بندوق} mean? (\textquotedblleft{}A gun\textquotedblright{}.) Say it once more.
\end{tcolorbox}
\section[\texorpdfstring{\ur{باندھنا} (bāndhnā)}{باندھنا (bāndhnā)}]{\ur{باندھنا} (bāndhnā) — to tie up; to bind}

\label{lesson:UR-C280-bandhna}



\begin{quote}
Before the new one: say the Urdu for an armed robbery, then the Urdu for a gun.
\end{quote}

\subsection*{You'll want to know: \ur{باندھنا}}
\textbf{\ur{باندھنا}} — \emph{bāndhnā} — \textquotedblleft{}to tie up; to bind\textquotedblright{}.

\begin{quote}
\ur{ڈاکوؤں} \ur{نے} \ur{چوکیدار} \ur{کو} \ur{باندھا۔}
\end{quote}

\begin{quote}
\emph{ḍākuoṅ ne chaukīdār ko bāndhā.}
\end{quote}

\begin{quote}
The bandits tied up the watchman.
\end{quote}

Past, with \ur{نے} (he or she …): \ur{اس} \ur{نے} \ur{رسی} \ur{باندھی}, \emph{us ne rassī bāndhī}.

\begin{cousinweb}
\textbf{\ur{باندھنا}} (\emph{bāndhnā}) is inherited from Sanskrit \emph{bandhati}, \textquotedblleft{}binds\textquotedblright{} (CDIAL 9139). A person tied up takes \textbf{\ur{کو}} and the verb stays masculine, \textbf{\ur{چوکیدار} \ur{کو} \ur{باندھا}}; a thing tied agrees, \textbf{\ur{رسی} \ur{باندھی}}.
\end{cousinweb}

\begin{grammarlens}[title={What you've built}]
One more word for telling what happened.
\end{grammarlens}

\subsection*{Your turn}
\begin{itemize}
\raggedright
  \item \emph{Say it:} \emph{bāndhnā}
  \item \emph{Say it:} \emph{bāndhnā}, once more
  \item \emph{Say it:} \emph{ḍākuoṅ ne chaukīdār ko bāndhā.} — The bandits tied up the watchman.
  \item \emph{Recall:} say the Urdu for an armed robbery, then the Urdu for a gun, then say \emph{bāndhnā} again
\end{itemize}

\begin{tcolorbox}[breakable,colback=teal!4,colframe=teal!35!black,title={Before you move on}]
What does \ur{باندھنا} mean? (\textquotedblleft{}To tie up; to bind\textquotedblright{}.) Say it once more.
\end{tcolorbox}
\section[\texorpdfstring{\ur{لٹیرا} (luṭerā)}{لٹیرا (luṭerā)}]{\ur{لٹیرا} (luṭerā) — a robber; a looter}

\label{lesson:UR-C280-lutera}



\begin{quote}
Before the new one: say the Urdu for a gun, then the Urdu for to tie up.
\end{quote}

\subsection*{You'll want to know: \ur{لٹیرا}}
\textbf{\ur{لٹیرا}} — \emph{luṭerā} — \textquotedblleft{}a robber; a looter\textquotedblright{}.

\begin{quote}
\ur{اگر} \ur{لٹیرے} \ur{آئے} \ur{تو} \ur{ہم} \ur{بھاگیں} \ur{گے۔}
\end{quote}

\begin{quote}
\emph{agar luṭere āe to ham bhāgeṅ ge.}
\end{quote}

\begin{quote}
If the robbers come, we will run.
\end{quote}

Masculine.

\begin{cousinweb}
\textbf{\ur{لٹیرا}} (\emph{luṭerā}) is built on \emph{lūṭnā}, \textquotedblleft{}to plunder\textquotedblright{}, which CDIAL lists for Hindi under \emph{*luṭṭati}, a form it reconstructs (CDIAL 11078), with the ending \emph{-erā} for the one who does it.
\end{cousinweb}

\begin{grammarlens}[title={What you've built}]
One more word for telling what happened.
\end{grammarlens}

\subsection*{Your turn}
\begin{itemize}
\raggedright
  \item \emph{Say it:} \emph{luṭerā}
  \item \emph{Say it:} \emph{luṭerā}, once more
  \item \emph{Say it:} \emph{agar luṭere āe to ham bhāgeṅ ge.} — If the robbers come, we will run.
  \item \emph{Recall:} say the Urdu for a gun, then the Urdu for to tie up, then say \emph{luṭerā} again
\end{itemize}

\begin{tcolorbox}[breakable,colback=teal!4,colframe=teal!35!black,title={Before you move on}]
What does \ur{لٹیرا} mean? (\textquotedblleft{}A robber; a looter\textquotedblright{}.) Say it once more.
\end{tcolorbox}
